\documentclass[10pt,a4paper]{amsart}
\usepackage[margin=19mm]{geometry}
\usepackage{amsmath,amssymb,mathtools}
\usepackage[scr=rsfs,cal=euler]{mathalfa}
\usepackage{xfrac}
\usepackage[unicode,hidelinks,bookmarks=false]{hyperref}
\newcommand{\ie}{i.e.\ }
\newcommand{\cf}{cf.\ }
\newcommand{\resp}{respectively\ }
\newcommand{\Subsection}{\subsection}
\providecommand{\nobreakdash}{\nobreak\hbox{-}}
\setlength{\emergencystretch}{2em}
\multlinegap0pt
\usepackage{bm}
\usepackage{bbm}
\usepackage{dsfont}
\usepackage{xspace}
\usepackage{cancel}
\usepackage{mathtools}
\usepackage{amsthm}
\makeatletter
\@ifundefined{theorem}{\newtheorem{theorem}{Theorem}}{}
\@ifundefined{definition}{\newtheorem{definition}[theorem]{Definition}}{}
\@ifundefined{lemma}{\newtheorem{lemma}[theorem]{Lemma}}{}
\makeatother
\makeatletter
\DeclareMathOperator{\Ad}{Ad}
\newcommand{\Thetahat}{\widehat{\Theta}}
\expandafter\ifx\csname authornote\endcsname\relax
\newenvironment{authornote}[1]{\addtocounter{anote}{1}\medskip \hrule \medskip \begin{quote}\small\noindent{\bf Note \arabic{anote}: #1} }
{\end{quote}\medskip \hrule\medskip}
\fi
\DeclareMathOperator{\curv}{curv}
\newcommand{\intSone}{\int_{0}^{2\pi}\mathrm{d}\theta\,}
\newcommand{\phihat}{\widehat{\phi}}
\DeclareMathOperator{\pr}{pr}
\makeatother
\title[Rigid models for 2-gerbes I: Chern-Simons geometry]{Rigid models for 2-gerbes I: Chern-Simons geometry}
\author{David Michael Roberts, Raymond F. Vozzo}
\date{Source: arXiv:2209.05521v4 (CC BY 4.0)}
\begin{document}
\maketitle

\noindent\emph{Source and licence.} Rigid models for 2-gerbes I: Chern-Simons geometry. Version arXiv:2209.05521v4 (CC BY 4.0), distributed under the Creative Commons Attribution 4.0 International licence, as recorded in the arXiv licence metadata for this exact version and shown on the abstract page https://arxiv.org/abs/2209.05521v4. Full rights notice: licenses/arxiv-2209.05521-CC-BY-4.0.txt.\par\smallskip

\noindent\textit{Editorial scope.} Counted unit: the Lemma whose source label is \texttt{lemma:rigid\_conn\_str\_2} in the article (source0.tex lines 2133--2135, source label \texttt{lemma:rigid\_conn\_str\_2}), delivered together with the article's own complete original proof of it (source0.tex lines 2137--2244, one \texttt{proof} environment of 108 lines). No mathematics is written, repaired or re-derived: statement and proof are the article's own text line for line. Required context retained uncounted: def:conn\_structure (lines 1199--1212). Every definition, display and label of those lines is retained unchanged. Omitted from this package and not redistributed: 1--1198, 1213--2132, 2136--2136, 2245--3325. Nothing inside a retained excerpt is omitted or altered: every retained body is delivered exactly as the article's lines are written, so no omission receipt of any kind accompanies this package. Citations: the retained excerpts carry no citation command at all, so no bibliography entry of the article is transcribed and no reference is redistributed. Reference adaptation: none was needed and none was made; every reference inside the delivered unit resolves inside it, so no unresolved reference or citation is produced. Printed slips kept literal: no printed word, symbol, hypothesis or number of the counted statement or of its proof was altered. Layout and class: the article's own class is \texttt{amsart} with packages including fontenc, amssymb, amsthm, mathrsfs, geometry, lscape, rotating, pdflscape, afterpage, xcolor, graphicx, comment, hyperref, enumerate; the wrapper is the lane's pinned \texttt{amsart} editorial wrapper at 10pt on A4, which restates the theorem-like environments the retained text needs and adds the article's own macro definitions that the retained text needs (\texttt{\textbackslash Ad}, \texttt{\textbackslash Thetahat}, \texttt{\textbackslash curv}, \texttt{\textbackslash intSone}, \texttt{\textbackslash phihat}, \texttt{\textbackslash pr}), with extra wrapper packages (bm, bbm, dsfont, xspace, cancel, mathtools). The delivered numbering is the unit's own. Assets: no retained span contains a figure environment, a graphics-inclusion command, a PDF-page inclusion, a table or a TikZ picture; no image file is copied into this package, the package declares no asset dependency and no image file is redistributed with it. Licence: version arXiv:2209.05521v4 of this article is distributed under the Creative Commons Attribution 4.0 International licence, as recorded in the arXiv licence metadata for this exact version and shown on the abstract page https://arxiv.org/abs/2209.05521v4. A full rights notice is delivered at licenses/arxiv-2209.05521-CC-BY-4.0.txt. Characters: every retained excerpt is pure ASCII, so the delivered engine, XeLaTeX with its default Latin Modern text font, reports no missing character.
\begin{definition}\label{def:conn_structure}
A \emph{connective structure} on a rigid bundle 2-gerbe consists of the differential form data\footnote{Of course, connections on the $U(1)$-bundle $E$ are only an affine subspace of $\Omega^1(E)$, but the conditions here cut down the space a lot more.} (using the notation as before) $(\nabla_E,\beta,\lambda) \in \Omega^1(E) \oplus \Omega^2(Z_1) \oplus  \Omega^3(Y)$ such that
\begin{enumerate}
\item $\nabla_E$ is a bundle gerbe connection: $0=\delta_v(\nabla_E) \in \Omega^1(\delta_vE)$;
\item $\beta$ is a curving for $\nabla_E$, namely $\delta_v(\beta) = \curv(\nabla_E)$;
\item $\lambda$ is a 2-curving for $\beta$, namely $\pi^*\delta_h(\lambda) = d\beta$.
\end{enumerate}
We say a connective structure is \emph{rigid} if there is a specified 1-form $\alpha$ on $Z_2$ such that
\begin{enumerate}\setcounter{enumi}{3}
\item $\delta_v(\alpha) = M^*\delta_h(\nabla_E)$;
\item $d\alpha = \delta_h\beta$;
\item $\delta_h(\alpha)=0$.
\end{enumerate}
\end{definition}

\begin{lemma}\label{lemma:rigid_conn_str_2}
The condition (5) in Definition \ref{def:conn_structure} holds for $\beta_A$ and $\alpha$, namely $d\alpha = \delta_h\beta_A$.
\end{lemma}

\begin{proof}
Let us first calculate $d\alpha$, taking into account the fact that $\alpha$ is independent of the $Q$-coordinate, that $PG^2$ is a Lie group, and that $\alpha$ is left invariant in the first coordinate.
Take a pair of tangent vectors $(pX_1,qY_1),(pX_2,qY_2) \in T_{(p,q)}PG^2$, and evaluate, where by an abuse of notation we are suppressing any mention of coordinates on $Q$:
\begin{align*}
&d\alpha_{(p,q)}\left((pX_1,qY_1),(pX_2,qY_2)\right) \\
& = \tfrac12 \Big\{
(pX_1,qY_1)\cdot \alpha(X_2,qY_2) - (pX_2,qY_2)\cdot \alpha(X_1,qY_1) 
- \alpha(p[X_1,X_2],q[Y_1,Y_2])
\Big\}.
\end{align*}
To continue, we need the derivative $(pX,qY)\cdot \alpha(pX',qY')$, which is
\begin{align*}
(pX,qY)\cdot \alpha(pX',qY') &=
\left.\frac{d}{dt} \left\{ 2\intSone \langle X',\partial_\theta[q(1+tY+O(t^2))](1-tY+O(t^2))q^{-1}\rangle\right\}\right|_{t=0}\\
& = \begin{multlined}[t]
2\intSone  \frac{d}{dt}\langle X',(\partial_\theta q)(1+tY+O(t^2))(1-tY+O(t^2))q^{-1}\rangle\big|_{t=0} \\
+
 \frac{d}{dt}\langle X',qt\partial_\theta(Y)(1-tY)q^{-1}+O(t^2)\rangle\big|_{t=0}  
\end{multlined}\\
& = 2\intSone  \frac{d}{dt}\langle X',(\partial_\theta q)(1+O(t^2))q^{-1}\rangle\big|_{t=0} 
+\langle X',q\partial_\theta(Y)q^{-1}\rangle \\
& = 2\intSone  \langle X',q\partial_\theta(Y)q^{-1}\rangle.
\end{align*}
Using this, we find
\[
d\alpha_{(p,q)}\left((pX_1,qY_1),(pX_2,qY_2)\right)=
2\intSone \langle X_2,\Ad_q\partial Y_1\rangle - \langle X_1 ,\Ad_q \partial Y_2\rangle - \langle [X_1,X_2],\phihat_q\rangle.
\]
Let us check some particular evaluations of $B$, noting that it is left-invariant:
\begin{align*}
  B(Y_1,\Ad_q^{-1}X_2) 
  % &= \intSone \langle Y_1,-q^{-1}\partial q\,q^{-1} X_2q + \Ad_{q^{-1}}\partial X_2 + q^{-1}X_2 \partial q\rangle\\
  &= \intSone \langle Y_1 ,\Ad_{q^{-1}}[X_2,\phihat_q]\rangle + \langle Y_1 ,\Ad_{q^{-1}}\partial X_2\rangle
\end{align*}
and 
\begin{align*}
B(\Ad_{q^{-1}}X_1,\Ad_{q^{-1}}X_2) 
& = \intSone \langle q^{-1}X_1q,\partial(q^{-1}X_2 q)\rangle\\
% & = \intSone {\color{red}\langle -X_1\partial q,q^{-1}X_2\rangle + \langle X_1,\partial X_2\rangle + \langle q^{-1} X_1 ,X_2 \partial q \rangle}\\
& = \intSone \langle X_1,[X_2,\phihat_q]\rangle + \langle X_1,\partial X_2\rangle.
\end{align*}
We also need to calculate $-\delta_h \pi^*\langle\pr_1^*A, \pr_2^* \Thetahat^G\rangle = -\pi^* \delta_h\langle\pr_1^*A, \pr_2^* \Thetahat^G\rangle$, and so we need
\begin{align*}
  \delta_h\langle\pr_1^*A, \pr_2^* \Thetahat^G\rangle_{(y,g,h)}&=\langle A_{yg},\Thetahat_h\rangle - \langle A_y,\Thetahat_{gh}\rangle 
  + \langle A_y,\Thetahat_g\rangle\\
  &=\langle \Ad_{g^{-1}}A_y + \Theta_g,\Thetahat_h \rangle
  - \langle A_y,\Thetahat_g+\Ad_g\Thetahat_{h}\rangle 
  + \langle A_y,\Thetahat_g\rangle\\
  &= \langle A_y,\Ad_g\Thetahat_h\rangle + \langle\Theta_g,\Thetahat_h \rangle
  - \langle A_y,\Thetahat_g+\Ad_g\Thetahat_{h}\rangle 
  + \langle A_y,\Thetahat_g\rangle\\
  & = \langle\Theta_g,\Thetahat_h \rangle = \kappa_{(g,h)}.
\end{align*}
And now we can calculate $\delta_h(B)_{(p,q)}$:
\begin{align*}
&\delta_h(B)_{(p,q)}\left((pX_1,qY_1),(pX_2,qY_2)\right) \\
& = B(X_1,X_2) - B(\Ad_{q^{-1}}X_1+Y_1,\Ad_{q^{-1}}X_2+Y_2) + B(Y_1,Y_2)\\
 & = B(X_1,X_2)- B(\Ad_{q^{-1}}X_1,\Ad_{q^{-1}}X_2)
- B( Y_1,\Ad_{q^{-1}}X_2)
- B(\Ad_{q^{-1}}X_1, Y_2)
\\
& = \begin{multlined}[t]
\frac12 \intSone \langle X_1,\partial X_2\rangle - \langle X_2,\partial X_1\rangle  
-\langle\Ad_{q^{-1}}X_1,\partial(q^{-1}X_2q)\rangle
+\langle\Ad_{q^{-1}}X_2,\partial(q^{-1}X_1q)\rangle\\
- \langle Y_1, \partial(q^{-1}X_2q)\rangle + 
\langle \Ad_{q^{-1}}X_2 , \partial Y_1\rangle 
-\langle \Ad_{q^{-1}}X_1 , \partial Y_2\rangle
+ \langle Y_1, \partial(q^{-1}X_2q)\rangle 
\end{multlined}\\
\\
&=\begin{multlined}[t]
\frac12 \intSone \langle X_1,\partial X_2\rangle - \langle X_2,\partial X_1\rangle - \langle X_1,\partial X_2\rangle - \langle X_1,[X_2,\phihat_q] \rangle\\
-\langle Y_1,\partial(\Ad_{q^{-1}}X_2)\rangle 
+ \langle \Ad_{q^{-1}}, \partial Y_1\rangle
+ \langle X_2,\partial X_1\rangle + \langle X_2,[X_1,\phihat_q]\rangle\\
-\langle\Ad_{q^{-1}}X_1,\partial Y_2\rangle
+ \langle Y_2,\partial(\Ad_{q^{-1}}X_1)\rangle
\end{multlined}
\\
&=\begin{multlined}[t]
\frac12 \intSone \langle X_2,\Ad_q\partial Y_1\rangle
- \langle X_1,\Ad_q Y_2\rangle
-2\langle [X_1,X_2,\phihat_q\rangle\\
-\langle Y_1,\partial(\Ad_{q^{-1}}X_2)\rangle
+ \langle Y_2,\partial(\Ad_{q^{-1}}X_1)\rangle
\end{multlined}
\\
&=\begin{multlined}[t]
\frac12\intSone \Big\{\langle X_2,\Ad_q\partial Y_1\rangle 
+ \langle \partial Y_1, \Ad_{q^{-1}} X_2\rangle
- \langle X_1,\Ad_q\partial Y_2\rangle
- \langle \partial Y_2,\Ad_{q^{-1}}X_1 \rangle \\
- 2\langle [X_1,X_2,\phihat_q\rangle\Big\} + \tfrac12 \left(-\langle X_2(2\pi),\Ad_{q(2\pi)}Y_1(2\pi)\rangle + \langle X_1(2\pi),\Ad_{q(2\pi)}Y_2(2\pi)\rangle \right)
\end{multlined}\\
&=\begin{multlined}[t]
\intSone \Big\{\langle X_2,\Ad_q\partial Y_1\rangle
- \langle X_1,\Ad_q\partial Y_2\rangle
- \langle [X_1,X_2,\phihat_q\rangle\Big\} + \pi^*\kappa_{(p,q)}\left((pX_1,qY_1),(pX_2,qY_2)\right)
\end{multlined}\\
&= d\alpha_{(p,q)}\left((pX_1,qY_1),(pX_2,qY_2)\right) + \pi^*\kappa_{(p,q)}\left((pX_1,qY_1),(pX_2,qY_2)\right).
\end{align*}
Thus
\[
\delta_h\beta_A= \delta_h B - \delta_h\pi^*\langle\pr_1^*A, \pr_2^* \Thetahat^G\rangle = d\alpha+ \pi^*\kappa -\pi^*\kappa =d\alpha,
\]
as we needed to show.
\end{proof}

\end{document}
